%%%%%%%%%%%%%%%%%%%%%%%%%%%%%%%%%%%%%%%%%
% Szablon pracy dyplomowej
% Wydział Informatyki 
% Zachodniopomorski Uniwersytet Technologiczny w Szczecinie
% autor Joanna Kołodziejczyk (jkolodziejczyk@zut.edu.pl)
% Bardzo wczesnym pierwowzorem szablonu był
% The Legrand Orange Book
% Version 5.0 (29/05/2025)2023)
%
% Modifications to LOB assigned by %JK
%%%%%%%%%%%%%%%%%%%%%%%%%%%%%%%%%%%%%%%%%
\chapter{Instrukcja instalacji i uruchomienia}
\label{chapter:dodatek_A}

W tym dodatku opisano procedurę instalacji i uruchomienia systemu ,,Household Manager'' w środowisku lokalnym.

\section{Wymagania wstępne}

Do poprawnego uruchomienia systemu wymagane jest zainstalowanie następujących narzędzi:
\begin{itemize}
    \item \textbf{Docker Desktop} (wersja 4.0 lub nowsza) - do konteneryzacji serwera i bazy danych.
    \item \textbf{Node.js} (wersja 18 LTS lub nowsza) - środowisko uruchomieniowe dla aplikacji mobilnej.
    \item \textbf{Git} - system kontroli wersji.
    \item \textbf{Aplikacja Expo Go} - zainstalowana na fizycznym urządzeniu mobilnym (Android/iOS) do testowania aplikacji.
\end{itemize}

\section{Uruchomienie serwera (Backend)}

1. Sklonuj repozytorium projektu:
\begin{lstlisting}[language=bash]
git clone https://github.com/Dragodui/diploma-server.git
cd diploma-server
\end{lstlisting}

2. Utwórz plik konfiguracyjny \texttt{.env} na podstawie przykładu:
\begin{lstlisting}[language=bash]
cp .env.example .env
\end{lstlisting}

3. Uzupełnij brakujące zmienne środowiskowe w pliku \texttt{.env} (klucze API do AWS S3, Google OAuth itp.). Dla celów testowych lokalnie można pozostawić wartości domyślne dla bazy danych.

4. Uruchom kontenery Docker:
\begin{lstlisting}[language=bash]
docker compose up -d --build
\end{lstlisting}

5. Serwer będzie dostępny pod adresem \texttt{http://localhost:8080}. Dokumentacja Swagger UI dostępna jest pod adresem \texttt{http://localhost:8080/swagger/index.html}.

\section{Uruchomienie aplikacji mobilnej (Frontend)}

1. Przejdź do katalogu aplikacji klienckiej:
\begin{lstlisting}[language=bash]
cd client
\end{lstlisting}

2. Zainstaluj zależności projektu:
\begin{lstlisting}[language=bash]
npm install
\end{lstlisting}

3. Uruchom serwer deweloperski Expo:
\begin{lstlisting}[language=bash]
npx expo start
\end{lstlisting}

4. Zeskanuj wygenerowany kod QR za pomocą aplikacji Expo Go na swoim telefonie (upewnij się, że telefon i komputer są w tej samej sieci Wi-Fi) lub naciśnij \texttt{a}, aby uruchomić emulator Androida (wymaga zainstalowanego Android Studio).

\section{Konfiguracja zmiennych środowiskowych}

Poniżej przedstawiono przykładową konfigurację pliku \texttt{.env} dla serwera:

\begin{lstlisting}[caption=Przykładowy plik .env]
PORT=8080
DB_HOST=postgres
DB_USER=postgres
DB_PASSWORD=secret
DB_NAME=household_db
JWT_SECRET=super_secret_key_change_me
AWS_REGION=eu-central-1
AWS_ACCESS_KEY_ID=AKIA...
AWS_SECRET_ACCESS_KEY=...
\end{lstlisting}
